\section*{الفصل \ref{ch:b3:organometallic-bonding} --- الكيمياء العضوية الفلزية: الترابط والربائط}

\begin{solution}{exo:b3:organometallic-bonding:1}
$\ce{Fe(CO)5}$: $8 + 10 = 18$. $\ce{Mn2(CO)10}$: $7 + 10 + 1 = 18$ لكل Mn. $\ce{CpMo(CO)3H}$: $6 + 5 + 6 + 1 =
18$. $\ce{[PtCl3(C2H4)]-}$: $10 + 3 + 2 + 1 = 16$. $\ce{Cr($\eta^6$-C6H6)(CO)3}$: $6 + 6 + 6 = 18$. $\ce{Cp2ZrCl2}$:
$4 + 10 + 2 = 16$.
\end{solution}

\begin{solution}{exo:b3:organometallic-bonding:2}
Fe(0) $d^8$؛ Mn(0) $d^7$؛ Mo(II) $d^4$؛ Pt(II) $d^8$؛ Cr(0) $d^6$؛ Zr(IV) $d^0$.
\end{solution}

\begin{solution}{exo:b3:organometallic-bonding:3}
$\ce{[Mn(CO)6]+} > \ce{Cr(CO)6} > \ce{[V(CO)6]-}$: فكلما كان الفلز أغنى بالإلكترونات، ازداد
تبرعه العكسي في $\pi^*$ وضعفت الروابط C--O.
\end{solution}

\begin{solution}{exo:b3:organometallic-bonding:4}
الإيثان ($\ce{Mn(CO)5}$ و$\ce{CpFe(CO)2}$ \omterm{def:b3:organometallic-bonding:isolobal}{متساويا الفصوص} مع \ce{CH3})؛ والإيثان مرة أخرى؛ والتيتراهيدران
$\ce{(CH)4}$ وقد استُبدلت بثلاث من مجموعاته CH الشظيةُ $\ce{Co(CO)3}$.
\end{solution}

\begin{solution}{exo:b3:organometallic-bonding:5}
$fac$ ($C_{3v}$): حزمتان (\babelsublr{A$_1$ + E}). $mer$ ($C_{2v}$): ثلاث حزم (\babelsublr{2A$_1$ + B$_1$}).
\end{solution}

\begin{solution}{exo:b3:organometallic-bonding:6}
الفوسفين الضخم يزحم الفلز: فالتفكك يخفف هذا الإجهاد، فيكون
أسرع، وتُفضَّل الوسائط المنخفضة التناسق.
\end{solution}

\begin{solution}{exo:b3:organometallic-bonding:7}
معقد الكروم \omterm{def:b3:organometallic-bonding:carbene}{كربين فيشر} (المستبدل OMe، وCr(0)، وربائط CO): يهاجم
الأمينُ كربونَ \omterm{def:b3:organometallic-bonding:carbene}{الكربين} المحب للإلكترونات ويحل محل OMe. أما ألكيليدين
التنتالوم \omterm{def:b3:organometallic-bonding:carbene}{فكربين شروك}: يهاجم كربونه المحب للنواة كربونيل
الكيتون، فيعطي ألكينًا ووحدة Ta=O، كما في تفاعل فيتيغ.
\end{solution}

\begin{solution}{exo:b3:organometallic-bonding:8}
في الترتيب المتناوب: لا تداخل $\delta$، فالإلكترونان $\delta$ غير رابطين: رتبة الرابطة 3. $\ce{[Re2Cl8]^{4-}}$:
$\sigma^2\pi^4\delta^2\delta^{\ast2}$، رتبة الرابطة 3.
\end{solution}

\begin{solution}{exo:b3:organometallic-bonding:9}
مسافة \babelsublr{M$\cdots$H} قصيرة (الحيود، والأفضل حيود النيوترونات)؛ وعدد موجي منخفض لتمدد C--H؛
وإشارة رنين مغناطيسي نووي في المجال المرتفع مع اقتران C--H أحادي الرابطة مخفَّض.
\end{solution}

\begin{solution}{exo:b3:organometallic-bonding:10}
16: المعقدات المربعة المستوية $d^8$ ($\ce{[PtCl4]^{2-}}$، حفاز ويلكنسون)، التي يكون مدار إلكترونها الثامن عشر
مرتفعًا جدًا؛ \omterm{def:b3:organometallic-bonding:metallocene}{والميتالوسينات} للفلزات المبكرة مثل $\ce{Cp2ZrCl2}$، المزدحمة أكثر من أن تستقبل
ربائط أخرى. 17: $\ce{V(CO)6}$، الذي لا يستطيع أن يتثنّى لأسباب فراغية. 19: الكوبالتوسين،
الذي يقع إلكترونه الإضافي في مدار مضاد للترابط (ويُفقد بسهولة).
\end{solution}

\begin{solution}{exo:b3:organometallic-bonding:11}
CO مستقبل $\pi$ قوي: يخفض المستوى $t_{2g}$ ويجعل $\Delta_{\mathrm o}$ كبيرة، أعلى كثيرًا من
طاقة الاقتران: $t_{2g}^6$، ديامغناطيسي. وتقع انتقالات \babelsublr{$d$--$d$}، عند $\Delta_{\mathrm o}$، في فوق البنفسجي،
ولا يُمتص أي ضوء مرئي.
\end{solution}

\begin{solution}{exo:b3:organometallic-bonding:12}
التبرع من $\pi$ والتبرع العكسي في $\pi^*$ كلاهما يخفض رتبة الرابطة C--C. ففي ملح
زايزه (Pt(II)، تبرع عكسي معتدل) تطول الرابطة قليلًا؛ أما الفلز المنخفض التكافؤ
الغني بالإلكترونات فيتبرع تبرعًا عكسيًا قويًا ويقترب المعقد من حد
\omterm{def:b3:organometallic-bonding:dcd}{حلقي البروبان الفلزي}، برابطة C--C قريبة من الرابطة الأحادية.
\end{solution}

\begin{solution}{pb:b3:organometallic-bonding:1}
\textbf{1.} $\ce{Fe^{2+}}$ ($d^6$) $+ 2 \times 6 = 18$.
\textbf{2.} $8 + 2 \times 5 = 18$.
\textbf{3.} $9 + 10 = 19$.
\textbf{4.} $10 + 10 = 20$.
\textbf{5.} $18 - 1 = 17$.
\textbf{6.} Fe(II) $d^6$؛ Co(II) $d^7$؛ Ni(II) $d^8$؛ Fe(III) $d^5$.
\textbf{7.} $a_{1g}$ ($d_{z^2}$)، $e_{1g}$ ($d_{xz}$، $d_{yz}$)، $e_{2g}$ ($d_{xy}$، $d_{x^2-y^2}$).
\textbf{8.} يتداخل الزوج $e_{1g}$ بقوة مع التركيب $e_{1g}$ الممتلئ للحلقتين: مدارات رابطة
معظمها على الحلقتين، ومدارات $e_{1g}^\ast$ مضادة للترابط معظمها على الفلز، مرتفعة
الطاقة.
\textbf{9.} $e_{2g} \approx a_{1g} < e_{1g}^\ast$.
\textbf{10.} $e_{2g}^4a_{1g}^2$: لا إلكترون منفرد.
\textbf{11.} إلكترون واحد في $e_{1g}^\ast$: إلكترون منفرد واحد، $1.73\,\mu_B$.
\textbf{12.} إلكترونان في المدارين $e_{1g}^\ast$، متوازيا السبين: إلكترونان منفردان، $\sqrt{8} =
2.83\,\mu_B$.
\textbf{13.} $e_{2g}^3a_{1g}^2$: إلكترون منفرد واحد. والثقب في $e_{2g}^3$ يعطي حالة منحلة مداريًا
(E) ذات \omterm{def:b3:complex-spectra-magnetism:moment}{إسهام مداري}.
\textbf{14.} إلكترونه التاسع عشر في مدار مضاد للترابط ويُنزع بسهولة،
فيعطي كاتيون الكوبالتوسينيوم ذا الثمانية عشر إلكترونًا.
\textbf{15.} الإلكترون المنزوع يأتي من مدار شبه غير رابط: فلا تكاد
البنية تتغير (طاقة إعادة تنظيم صغيرة، \cref{ch:b3:complex-mechanisms}).
\textbf{16.} المزدوجة سريعة وعكوسة في معظم المذيبات، والمركب
ذواب ومستقر، وجهده قليل التأثر بالمذيب لأن
الشحنة موزعة على أيون كبير.
\textbf{17.} بعشرين إلكترونًا، اثنان منها مضادان للترابط، يتفاعل بفقدهما: فهو
يضيف ربائط مع انزلاق حلقة أو فقدها، أو يتأكسد.
\textbf{18.} $\ce{CpFe(CO)2}$: $8 + 5 + 4 = 17$ إلكترونًا، ينقصه واحد ليبلغ 18، ومدار حدودي واحد: فهو \omterm{def:b3:organometallic-bonding:isolobal}{متساوي الفصوص}
مع \ce{CH3}. والثنائي $\ce{[CpFe(CO)2]2}$ «إيثان»، برابطة Fe--Fe (وفي الواقع تجسر بينهما ربيطتا
CO).
\textbf{19.} $\ce{CpFe(CO)2-Mn(CO)5}$، برابطة Fe--Mn.
\textbf{20.} $\ce{PPh3}$ مانح أفضل ومستقبل $\pi$ أضعف من CO: فيزداد تبرع الفلز العكسي
إلى جزيئات CO الباقية، فتنخفض أعداد تمددها الموجية.
\textbf{21.} حزمتان (\babelsublr{A$_1$ + E}).
\textbf{22.} انزلاق الحلقة إلى $\eta^3$ يحرر ما يعادل إلكترونين من التناسق: فتستطيع ربيطة
داخلة أن ترتبط ترابطيًا دون أن يتجاوز الفلز 18 إلكترونًا.
\textbf{23.} \textbf{$\mu_{\text{السبين وحده}}(\text{النيكلوسين}) = \sqrt{2 \times 4} = 2.83\,\mu_B$.}
\end{solution}
